% ===============================================
% CONCLUSION ET PERSPECTIVES
% ===============================================

\chapter{Conclusion et perspectives} \label{ch_six}

\section{Bilan}

Ce mémoire a présenté la conception, l'implémentation et l'évaluation d'un pipeline modulaire de reconstruction 3D des mouvements de force athlétique à partir de vidéos monoculaires non calibrées. L'objectif central était de proposer une méthode accessible, sans marqueurs physiques ni d'infrastructure multicaméras, capable de gérer explicitement les occlusions générées par l'équipement de force athlétique.

Quatre contributions principales ont été développées au cours de ce travail.

Premièrement, un pipeline séquentiel opérationnel intégrant quatre modèles de l'état de l'art (SAM3, MoGe-2, TACO, SAM 3D Body) a été conçu. Cette architecture permet de passer d'une vidéo brute filmée en salle de musculation à un maillage 3D complet du corps de l'athlète, incluant le corps, les mains et les pieds, sans marqueurs physiques ni calibration préalable de la caméra. Le découpage modulaire facilite le remplacement de chaque composant au fil des avancées de l'état de l'art.

Deuxièmement, un dataset spécialisé pour la force athlétique a été constitué, regroupant des séquences comprenant les trois mouvements de référence (squat, développé couché, soulevé de terre) issues de six sujets distincts. Les occlusions synthétiques ont été générées via une banque d'occlusions spécialisée (disques IPF, disques de salle commerciale, racks à squat), et huit augmentations visuelles ont été appliquées sur les vidéos sources pour pallier la taille restreinte du corpus. La méthodologie de construction, entièrement automatisée et réutilisant les premiers modules du pipeline d'inférence, garantit la cohérence entre les distributions d'entraînement et d'inférence.

Troisièmement, un fine-tuning LoRA de TACO sur ce dataset a été conduit et évalué quantitativement. Sur un jeu de données de test composé de trois séquences d'un sujet absent du dataset d'entraînement, le modèle fine-tuné améliore le LPIPS de 35,0\,\% en moyenne par rapport au modèle baseline ($0{,}0577 \rightarrow 0{,}0375$), avec une significativité statistique confirmée par un test de Wilcoxon ($p = 3{,}28 \times 10^{-5}$). L'amélioration est particulièrement marquée au développé couché ($-50{,}1\,\%$ de LPIPS), mouvement pour lequel l'occlusion quasi continue du torse constitue précisément le cas d'usage ciblé par le fine-tuning. L'évaluation via le MPJPE confirme que cette amélioration a un impact en aval sur la reconstruction 3D : au développé couché, TACO-LoRA atteint un MPJPE de 37,2~mm, légèrement inférieur à celui obtenu sur l'image occluse brute (38,0~mm), et au soulevé de terre, le fine-tuning réduit l'erreur de 40,6\,\% par rapport au baseline (108,0~mm contre 181,8~mm).

Quatrièmement, une interface web de visualisation a été développée, permettant d'explorer les maillages 3D produits par SAM 3D Body et de calculer automatiquement les indicateurs biomécaniques pertinents pour chaque mouvement : profondeur du squat, verrouillage des coudes et des genoux, inclinaison du buste, et asymétrie des hanches. Cette interface constitue le premier maillon applicatif concret du pipeline.

\section{Limites} \label{sec:limites}

Plusieurs limites du pipeline développé doivent être reconnues.

\textbf{Taille modeste du jeu de données de test.} L'évaluation quantitative repose sur trois séquences issues d'un unique sujet, couvrant une diversité restreinte de morphologies, d'éclairages et de configurations de salle. Les résultats doivent donc être considérés comme indicatifs plutôt que définitifs, et leur généralisation à d'autres sujets et conditions de capture reste à valider sur un corpus plus large.

\textbf{Distribution d'occlusions identique entre entraînement et test.} Les séquences du jeu de données de test ont été occluses en utilisant la même banque d'occlusions synthétiques (disques IPF, racks) et le même pipeline de génération que les séquences d'entraînement. Cette homogénéité dans la distribution des occlusions favorise mécaniquement le modèle fine-tuné, qui a été spécifiquement entraîné sur ce type d'artefact. La performance sur des occlusions réelles de compétition, dont l'apparence (perspective, ombres portées, reflets) diffère nécessairement des occlusions synthétiques, reste donc à valider.

\textbf{Évaluation perceptive humaine absente.} Aucune étude utilisateur n'a été menée auprès d'un panel d'entraîneurs ou d'arbitres afin d'évaluer la qualité perçue des reconstructions. La validation repose exclusivement sur des métriques automatiques (LPIPS, SSIM, MPJPE). LPIPS corrèle mieux que SSIM avec le jugement humain mais ne le remplace pas, des reconstructions anatomiquement plausibles peuvent obtenir un bon LPIPS tout en présentant des artefacts subtils détectables à l'œil.

\textbf{Absence de comparaison avec un fine-tuning complet.} Le coût computationnel d'un fine-tuning complet de TACO étant prohibitif sur le matériel disponible ($2 \times$ RTX A5000 contre $8 \times$ H100 requis), l'écart de performance entre l'adaptation LoRA et un fine-tuning complet ne peut être quantifié dans ce travail. Il est raisonnable de supposer qu'un fine-tuning complet apporterait un gain supplémentaire, au prix d'un risque accru de catastrophic forgetting sur un dataset de cette taille.

\textbf{Artefacts de complétion amodale perturbant la reconstruction 3D.} Pour deux des trois mouvements évalués (squat et soulevé de terre), TACO-LoRA produit un MPJPE supérieur à celui obtenu directement sur l'image occluse brute (+3,6~mm et +13,7~mm respectivement). Ce résultat contre-intuitif indique que la complétion amodale, même améliorée par le fine-tuning, peut introduire des artefacts visuels (silhouettes légèrement déformées, textures incohérentes aux frontières de l'occlusion) suffisamment importants pour perturber SAM 3D Body davantage que l'occlusion originale. Ce phénomène suggère que la qualité perceptuelle mesurée par LPIPS et la qualité géométrique mesurée par MPJPE ne sont pas parfaitement corrélées, et qu'une amélioration de la complétion amodale 2D ne se traduit pas nécessairement par une amélioration de la reconstruction 3D en aval.

\textbf{Absence de contexte temporel global pour TACO.} Le traitement par fenêtres de 14 trames prive TACO du contexte global du mouvement. Les frontières entre fenêtres consécutives peuvent présenter de légères discontinuités visuelles, ce qui peut limiter la qualité des reconstructions sur les trames proches de ces frontières.

\textbf{Absence de cohérence temporelle 3D.} SAM 3D Body opère trame par trame sans mécanisme de lissage temporel des paramètres articulaires. Les reconstructions consécutives peuvent présenter des discontinuités dans les angles articulaires estimés, ce qui limite la qualité des analyses biomécaniques dérivées (vitesses angulaires, détection des \textit{sticking points}).

\textbf{Évaluation limitée aux séquences d'entraînement.} Le pipeline n'a pas été testé sur des séquences de compétition officielle, dont les conditions (éclairage de scène, public en arrière-plan, multiples points de vue) diffèrent significativement des salles d'entraînement.

\section{Travaux futurs}

Plusieurs axes d'amélioration et d'extension sont identifiés.

\textbf{Traitement temporel global de TACO.} L'approche idéale consisterait à traiter la séquence entière en une seule passe, en fournissant à TACO l'intégralité des trames comme contexte temporel. À court terme, l'introduction d'un chevauchement à 50\,\% entre fenêtres consécutives avec fusion par interpolation constituerait une amélioration immédiate, au prix d'un doublement du temps d'inférence.

\textbf{Évaluation du catastrophic forgetting.} Une évaluation sur un sous-ensemble du dataset OvO original permettrait de quantifier la dégradation des capacités généralistes du modèle suite au fine-tuning spécialisé, et d'ajuster le rang LoRA ou le taux d'apprentissage en conséquence.

\textbf{Comparaison avec les alternatives.} Une évaluation formelle de pix2gestalt et ProPainter sur le jeu de données de test powerlifting permettrait de valider quantitativement les arguments avancés au chapitre~\ref{ch_three} en faveur de TACO.

\textbf{Fine-tuning de SAM 3D Body.} Les résultats MPJPE montrent que SAM 3D Body produit des erreurs élevées sur le soulevé de terre, suggérant que le modèle généraliste peine sur les poses extrêmes de ce mouvement. Un fine-tuning spécialisé sur des séquences de force athlétique annotées en 3D constituerait une amélioration. Par ailleurs, comme discuté en section~\ref{sec:limites}, la complétion amodale par TACO introduit parfois des artefacts qui dégradent la reconstruction 3D, adapter SAM 3D Body aux occlusions caractéristiques du domaine permettrait de relâcher la dépendance à une complétion amodale 2D parfaite, et de bénéficier directement de la pose réelle même en présence d'occlusion partielle.

\textbf{Cohérence temporelle 3D.} L'intégration d'un lissage temporel des paramètres MHR entre trames consécutives (par exemple via un filtre de Kalman sur les angles articulaires) améliorerait la qualité des analyses biomécaniques et réduirait les discontinuités articulaires.

\textbf{Extension au domaine compétition.} La collecte de séquences filmées en conditions de compétition officielle (éclairage de scène, présence du public, prises de vue multiples) permettrait d'évaluer et d'améliorer la robustesse du pipeline dans ces conditions plus exigeantes, ouvrant la voie à une utilisation comme outil d'assistance à l'arbitrage.

\textbf{Optimisation et déploiement en temps réel.} Le pipeline actuel repose sur quatre modèles fondationnels de grande taille, dont certains comptent plusieurs milliards de paramètres. Cette architecture tire parti de la généralisation \textit{zero-shot} offerte par ces modèles, au prix d'un temps d'inférence incompatible avec un traitement en temps réel sur du matériel grand public. Des travaux futurs pourraient explorer des techniques de compression de modèle pour réduire ce coût sans sacrifier les capacités de généralisation : la distillation de connaissance (\emph{knowledge distillation})~\cite{Distilling-2015} permettrait d'entraîner des modèles étudiants plus légers en reproduisant le comportement des modèles enseignants sur le domaine de la force athlétique, la quantification (\emph{quantization})~\cite{Quantification-2021} réduirait l'empreinte mémoire en passant de la précision flottante 32 bits à des représentations 8 ou 4 bits, et l'élagage (\emph{pruning})~\cite{Brain-Damage-1989} supprimerait les poids peu contribuants. Ces techniques pourraient ouvrir la voie à un déploiement temps réel sur du matériel embarqué ou mobile, rendant l'outil accessible directement en salle d'entraînement sans infrastructure GPU dédiée.

\textbf{Amélioration des métriques biomécaniques.} Les indicateurs calculés dans l'interface actuelle constituent une première approximation fondée sur un sous-ensemble limité de keypoints. Des travaux futurs pourraient enrichir cette analyse en exploitant l'intégralité des 70 keypoints produits par SAM 3D Body (corps, mains et pieds), notamment pour calculer les angles articulaires des chevilles (indicateur de mobilité au squat), la trajectoire de la barre reconstruite depuis les poignets, la vitesse angulaire des articulations clés, ou encore la détection automatique des \emph{sticking points}. L'intégration des paramètres de forme MHR permettrait également d'estimer les moments de force articulaires, ouvrant la voie à une analyse biomécanique quantitative comparable à celle obtenue par les systèmes marker-based en laboratoire.

\textbf{Élargissement de la diversité du dataset de fine-tuning.} Le dataset constitué dans ce travail ne couvre que six sujets distincts, limitant la diversité morphologique et vestimentaire représentée durant l'entraînement. Une collecte étendue à plusieurs dizaines de sujets, intégrant explicitement les axes de variabilité décrits au chapitre~\ref{ch_four} (différents sexes, gabarits, tenues de compétition, conditions d'éclairage), permettrait d'améliorer significativement la généralisation du modèle aux populations non vues lors de l'entraînement.

\textbf{Amélioration de l'algorithme d'occlusion synthétique.} La banque d'occlusions actuelle se limite à des images statiques de disques et de racks superposés selon une trajectoire verticale simple. Une génération d'occlusions plus réaliste pourrait intégrer des modèles 3D rendus dynamiquement avec gestion des ombres portées, des reflets sur les surfaces métalliques, et de la déformation de perspective selon l'angle de prise de vue. Une telle approche réduirait l'écart de domaine entre occlusions synthétiques et occlusions réelles de compétition, améliorant directement la transposabilité du fine-tuning aux conditions in-the-wild.

\textbf{Sélection du sujet d'intérêt en amont du traitement.} Le pipeline actuel isole le sujet d'intérêt après segmentation par SAM3 et MoGe-2 via un filtre de profondeur. Une alternative consisterait à identifier le sujet d'intérêt en amont (par exemple via un prompt visuel sur la première trame), ce qui permettrait d'exploiter directement le mécanisme de suivi temporel intégré à SAM3. Cette approche garantirait la continuité d'identité du sujet entre trames consécutives, supprimerait le besoin du filtre de profondeur dans la majorité des cas, et améliorerait la robustesse du pipeline aux scènes contenant plusieurs personnes proches visuellement.